\begin{abstract}
Developers routinely choose open-source licenses, or inherit them through dependencies, without fully understanding the obligations involved. Large language models (LLMs) can explain legal text in plain language, but they also produce fluent, unsupported statements, which is especially harmful in licensing. We present \sys{}, a two-stage LLM\,+\,NLI pipeline that takes a natural-language project description and/or a dependency manifest, recommends three licenses from a catalogue of 34 OSI-approved software licenses, and explains each recommendation with plain-language claims that cite and quote a specific license clause. An off-the-shelf natural language inference (NLI) model then checks every claim against its cited clause and labels it \emph{verified}, \emph{unverified}, or \emph{contradicted}. Because no legal domain expert was available, we build all ground truth from public, machine-readable sources: choosealicense.com labels and texts, the SPDX license list, and the Software Heritage license dataset. We evaluate recommendation on 104 model-written, round-trip-validated project descriptions covering 70 requirement profiles, and verification on 729 clause--claim pairs, 146 of which were hand-checked by two annotators (Cohen's $\kappa = 0.88$). The recommender (gpt-oss-120b) places a correct license among its top three for 74\% of descriptions, falling to 60\% when only one to three licenses satisfy the requirements. The verifier (DeBERTa-v3-large) reaches 0.82 accuracy with 0.94 precision on \emph{verified}, but confirms only 69\% of true claims, struggling most with copyleft and commercial-use claims. These results suggest that LLMs are useful for interpreting developer intent and explaining licenses, but that license selection itself should be constrained by structured license knowledge.
\end{abstract}
